\documentclass[a4paper,12pt]{article}
\usepackage[utf8]{inputenc}
\usepackage[T2A]{fontenc}
\usepackage[russian]{babel}
\usepackage{amsmath,amssymb,amsthm}
\usepackage{mathtools}

\newtheorem{theorem}{Теорема}
\newtheorem{lemma}{Лемма}
\newtheorem*{problem}{Задача}
\renewcommand{\proofname}{Решение}

\begin{document}

\begin{problem}[5]
\end{problem}
\begin{proof}
$\implies$ Очевидно: если $f(z)=\sum_{j=0}^\infty x_j(z-z_0)^j$ в смысле сходимости по норме, то для любого $x^*\in X^*$ по непрерывности и линейности функционала
\[
x^*(f(z))=\sum_{j=0}^\infty x^*(x_j)(z-z_0)^j
\]
--- аналитическая функция, а т.к. $z_0$ произвольно, то $\varphi_{x^*}$ целая $\implies f$ --- слабо аналитична.

$\impliedby$ Докажем в обратную сторону

Зафиксируем $z_0\in\mathbb{C}, \> r>0$, $\gamma$ --- окружность $|\zeta-z_0|=r$. Для $|z-z_0|<r$ определим
\[
g(z)=\frac{1}{2\pi i}\int_\gamma \frac{f(\zeta)}{\zeta-z}\,d\zeta
\]
Интеграл понимается как интеграл Римана от непрерывной функции $\zeta\mapsto f(\zeta)/(\zeta-z)$ со значениями в $X$. Обоснование почему так можно в конце. 
Для любого $x^*\in X^*$ в силу непрерывности и линейности
\[
x^*(g(z))=\frac{1}{2\pi i}\int_\gamma \frac{x^*(f(\zeta))}{\zeta-z}\,d\zeta
\]
Внутренний интеграл --- интегральная формула Коши для $\varphi_{x^*}$, поэтому $x^*(g(z))=\varphi_{x^*}(z)=x^*(f(z))$. 
Поскольку $X^*$ разделяет точки $X$, то $g(z)=f(z)$ для всех $|z-z_0|<r$. Таким образом,
\[
f(z)=\frac{1}{2\pi i}\int_\gamma \frac{f(\zeta)}{\zeta-z}\,d\zeta \qquad (|z-z_0|<r)
\]

Теперь, раз у нас есть формула Коши, то можно провести аналогичные с классическими рассуждения и получить формулы для коэффициентов.
Для $|z-z_0|<r$ и $\zeta\in\gamma$ имеем:
\[
\frac{1}{\zeta-z}=\frac{1}{\zeta-z_0}\cdot\frac{1}{1-\frac{z-z_0}{\zeta-z_0}}
= \sum_{j=0}^\infty \frac{(z-z_0)^j}{(\zeta-z_0)^{j+1}}
\]
Т.к. $\|f(\zeta)\|\le M$ на $\gamma$, то 
\[
\sum_{j=0}^\infty \frac{f(\zeta)(z-z_0)^j}{(\zeta-z_0)^{j+1}}
\]
сходится равномерно по $\zeta\in\gamma$ в смысле нормы $X$, тогда
\[
f(z)=\frac{1}{2\pi i}\int_\gamma \sum_{j=0}^\infty \frac{f(\zeta)(z-z_0)^j}{(\zeta-z_0)^{j+1}}\,d\zeta
= \sum_{j=0}^\infty \left(\frac{1}{2\pi i}\int_\gamma \frac{f(\zeta)}{(\zeta-z_0)^{j+1}}\,d\zeta\right) (z-z_0)^j
\]
Наконец положим
\[
x_j:=\frac{1}{2\pi i}\int_\gamma \frac{f(\zeta)}{(\zeta-z_0)^{j+1}}\,d\zeta \ \in X.
\]
Тогда
\[
f(z)=\sum_{j=0}^\infty x_j(z-z_0)^j
\]
для $|z-z_0|<r$, причём ряд сходится по норме. Поскольку $z_0$ и $r$ произвольны, то мы получаем сильную аналитичность $f$ на $\mathbb{C}$.
\end{proof}

\paragraph{\textbf{Почему можно интегрировать}}

Зафиксируем произвольный круг $D=\{|z-z_0|\le r\}$. $ \forall x^*\in X^*$ функция $\varphi_{x^*}(z)=x^*(f(z))$ непрерывна, тогда
$\sup_{z\in D}|\varphi_{x^*}(z)|<\infty$


Рассмотрим семейство $\{f(z):z\in D\}\subset X$. Можно представлять, что они действуют на $X^*$ как элементы $X^{**}$ соответствующим образом:
$f(z)(x^*)=x^*(f(z))$. Т.к. для всякого $x^*$ множество $x^*(f(D))$ ограничено, то семейство $\{f(z)\}$ поточечно ограничено на $X^*$. 
И поскольку $X^*$ --- банахово, то по принципу равномерной ограниченности оно ограничено по норме в $X^{**}$:
\[
\sup_{z\in D}\|f(z)\|_{X^{**}}<\infty
\]
Т.к. $X$ полное, то из ограниченности в $X^{**}$ следует ограниченность в $X$:
\[
\|f(z)\|_X=\|in(f(z))\|_{X^{**}}\le \sup_{w\in D}\|f(w)\|_{X^{**}}<\infty.
\]
где $in: X \to X^{**}$ --- каноническое вложение. Итого, $\exists M>0$, т.ч. $\|f(z)\|\le M, \> \forall z\in D$.

Далее, рассмотрим тот же круг $D$, и $z,w \in D$, т.ч. $|z-z_0|<r/2$, $|w-z_0|<r/2$, тогда $\forall x^*\in X^*$ по формуле Коши имеем:
\[
\varphi_{x^*}(z)-\varphi_{x^*}(w)=\frac{1}{2\pi i}\int_\gamma \varphi_{x^*}(\zeta)\left(\frac{1}{\zeta-z}-\frac{1}{\zeta-w}\right)d\zeta
= \frac{z-w}{2\pi i}\int_\gamma \frac{\varphi_{x^*}(\zeta)}{(\zeta-z)(\zeta-w)}\,d\zeta.
\]
Учитывая, что $\varphi_{x^*}(\zeta)=x^*(f(\zeta))$, получаем:
\[
x^*(f(z)-f(w))=\frac{z-w}{2\pi i}\int_\gamma \frac{x^*(f(\zeta))}{(\zeta-z)(\zeta-w)}\,d\zeta.
\]
Оценим модуль. Для $\zeta\in\gamma$ выполнено $|\zeta-z|\ge r/2$, $|\zeta-w|\ge r/2$ и $\|f(\zeta)\|\le M$. Тогда
\[
|x^*(f(z)-f(w))|\le \frac{|z-w|}{2\pi}\cdot 2\pi r\cdot \frac{M\|x^*\|}{(r/2)^2}= \frac{4M}{r}|z-w|\,\|x^*\|.
\]
Итого, $\forall \|x^*\|\le 1$ имеем
\[
\left| x^*\left(\frac{f(z)-f(w)}{z-w}\right) \right|\le \frac{4M}{r}.
\]
Значит, множество $\left\{\frac{f(z)-f(w)}{z-w}:\ |z-z_0|,|w-z_0|<r/2\right\}$ слабо ограничено. Снова по принципу равномерной ограниченности 
оно ограничено по норме:
\[
\|f(z)-f(w)\|\le K|z-w|,
\]
где $K=4M/r$. Т.е. $f$ локально липшицева.

После этого вроде как должны нормально работать все рассуждения с суммами Дарбу и проч. Т.е. можно нормально определить интеграл Римана


\end{document}

